\documentclass{article}
\usepackage[utf8]{inputenc}
\usepackage[T1]{fontenc}
\usepackage{geometry}
\usepackage{graphicx}
\usepackage{float}
\usepackage{hyperref}
\usepackage[english,indonesian]{babel}

\title{Optimasi Pemrosesan Template LaTeX Menggunakan Pendekatan Berbasis AI}
\author{Academic Support Agent}
\date{\today}

\begin{document}
\maketitle

\begin{abstract}

\end{abstract}

\chapter{PENDAHULUAN}


\chapter{TINJAUAN PUSTAKA}


% Original content (commented out):
% %  +-------------------------------------+
%  | LaTeX class aesbr.cls template file |
%  +-------------------------------------+
% To be used with "aesbr.cls" version 1.0b

%%
%% NOTE: This text file uses UNIX line feed conventions. When (human)
%% reading this file on other platforms, you may have to use a text
%% editor that can handle lines terminated by the UNIX line feed
%% character (0x0A).
%%

%%%%%%%%%%%%%%%%%%%%%%%%%%%%%%%%%%%%%%%%%%%%%%%%%%%%%%%%%%%
%
%  LOADING THE CLASS
%
% The class "aesbr.cls" automatically loads the following packages:
%    hyperref, graphicx, color, times, and helvet
% The package "graphicx.sty" are required if you want graphics,
% photos, etc. Be warned that pdflatex will require
% graphics to be in PDF, PNG, or JPG (not EPS) format.

% Para gerar documento no formato DVI, remova a opção 'pdfout' execute:
% > latex template
% > bibtex template
% > latex template
%
% Para gerar o documento no formato PDF execute:
% > pdflatex template
% > bibtex template
% > pdflatex template

% Remove the option "pdfout" if you are not using 'pdflatex'.
% The option "ams" loads the packages "amsmath", "amssymb" and "amsthm".
\documentclass
  [ams,pdfout]% class options
	{aesbr}

%%%%%%%%%%%%%%%%%%%%%%%%%%%%%%%%%%%%%%%%%%%%%%%%%%%%%%%%%%%
%
% LOADING PACKAGES
%
% Use \RequirePackage instead of \usepackage.
% It prevents a package clash.

% Use 8-bit UCS/Unicode Transformation Format
\RequirePackage[utf8]{inputenc}

% Para escrever em portugues é necessário usar a opção "brazil" do "babel"
% If you are writing in english change the order of the language options.
% You should use the "brazil" option even if you are writing in english.
\RequirePackage[english,brazil]{babel}
%\RequirePackage[brazil,english]{babel}

% Caracteres extendidos
% Extended characteres
\RequirePackage{textcomp}
\RequirePackage{url}

\begin{document}
% A página de título é construida no ambiente "TitlePage".
% Nesse ambiente deve-se informar, pela ordem, o título
% do artigo, os autores e respectivas afiliações, e o
% resumo.
% The title page is composed with the "TitlePage" environment.
% The information passed to this environment should follow the
% following order: title, authors with affiliations, and the abstract.
\begin{TitlePage}
	% O título é produzido com o comando \Title.
	% The title is produced with the command \Title
	\Title{Título do Artigo}
	% \RunningTitle e \RunningAuthors serão impressos
	% no topo de cada página
	% \RunningTitle and \RunningAuthors will be printed
	% on the head of each page
	\RunningTitle{Título Curto do Artigo} % short title
	% 2 or more authors:
    % ("Author1 and Author2" for 2 authors.)
    % ("Author1 et al." for 3 or more authors.)
	\RunningAuthors{Fulano et al.}
% Specify \Author and \Affil for each author (in this order)
% First author
	\Author
		[fulano@email.org]% email
		{Fulano de Tal}% Author name (printed in title page)
		\Affil
			[A]% chave/key (optional argument)
			{%
				Universidade,
				Instituto,
				Departamento\\
				Cidade, Estado, Código Postal, País
			}
		% Se uma mesma afiliação for compartilhada
		% por um ou mais autores, use uma <chave> (A, neste caso)
		% (qualquer combinação de caracteres do alfabeto).
    % Para se referir a esta afiliação
		% posteriormente, use o comando \Affilref[<chave>].
		%
		% If the same affiliation is for more than one author
		% you can use the second optional parameter <key>
		% (any character combination in ASCII) as a label. This label
    % can then be referred when you specify another author with
    % the same affiliation. To refer to that shared affiliation
    % use the command \Affilref[<key>]
% Second Author
	\Author
		[cicrano@email.com]
		{Cicrano de Tal}
		\Affilref[A]% mesma afiliação de "fulano" (chave A)
	                           % the same affiliation as "fulano" (key A)
% Third author (share the affiliation with Fulano de Tal)
	\Author
		[beltrano@email.fr]
		{Beltrano de Tal}
		\Affil
			{%
				University,
				Department of Civil Engineering\\
				City, State, Postal Code, Country
			}%
	\Abstract{%

O artigo deve conter um resumo ou \emph{abstract} ou \emph{resumen} de  60 a 120 palavras. Todos os manuscritos devem ser preparados em coluna dupla e folha formato A4, não numeradas. Cite os autores adequadamente, indicando sua afiliação. Certifique-se que um “e” seja colocado antes do último autor. Para cada afiliação indique seu endereço, cidade, estado, código postal e país. Colocar sempre as afiliações uma embaixo da outra. O título abreviado do artigo (no topo à direita  da próxima página) deve ocupar somente uma linha, com alinhamento à direita, junto com o(s) sobrenome(s) do(s) autor(es) alinhado(s) à esquerda. O autor deverá adotar RESUMO para os artigos em português, ABSTRACT para os artigos em inglês ou RESUMEN para os artigos em espanhol.
	}%
\end{TitlePage}
%
% Digite os títulos das seções EM LETRAS MAIÚSCULAS
% Type the section titles IN CAPITAL LETTERS.
\section{SEÇÃO}
Este \emph{template,} em \LaTeX\, deve ser compatível com as principais distribuições do \LaTeX. O objetivo deste template é sugerir um formato padrão para apresentação de trabalhos técnicos e científicos. Para isto, basta salvar este template com outro nome do arquivo (por exemplo: \emph{Sobrenom\_AESBR2016.tex}), digitar o novo texto sobre este texto, e deletar aquilo que não for necessário.

Os artigos submetidos ao Congresso da AES não são revisados pelo corpo editor, e são publicados em suas formas originais, como submetidos. As versões finais devem ser enviadas preferencialmente em formato PDF (\emph{Sobrenome\_AESBR2016.pdf}).

Os titulos das seções devem ser escritos em letras maiúsculas. A seção 0 é usualmente uma INTRODUÇÃO, seguindo-se a convenção da AES.
% ou \emph{postscript} (*.ps).
%
\subsection{Sub-seção 1}
Sub-seções usam letras maiúsculas e minúsculas, como acima. Para produzir sub-seções sem numeração use a forma com asterisco: \verb|\subsection*|.

\section{CONTEÚDO}
Para garantir que os artigos do Congresso da AES sejam consistentes com os objetivos da AES Brasil, as instruções abaixo devem ser consideradas pelos autores.

O conteúdo técnico deve ser preciso e coerente. Citação a trabalhos anteriores e/ou de terceiros devem ter seus respectivos créditos.

Se o artigo descrever um produto, o conteúdo deve enfocar os aspectos técnicos deste produto (circuito, leiaute, especificações, funções, aplicações, etc.).

Logomarcas de empresas não devem ser usadas. Nomes de empresas e modelos de equipamentos não devem ser colocados no título ou no resumo, e devem ser mínimos no texto (devem-se usar descrições genéricas).

Símbolos e marcas registradas não podem ser colocados no título e/ou no resumo, e preferencialmente não devem ser usadas no texto. Caso marcas registradas apareçam no texto, estas devem indicar junto ao nome o símbolo ``\texttrademark'' e deve-se incluir uma nota de rodapé dizendo quem é (são) o(s) detentor(es) destas marcas. Isso deve ser feito na primeira vez que a(s) marca(s) registrada(s) aparecer(em).
%
\subsection{Referências Bibliográficas}

As referências devem ser numeradas e listadas ao final do texto, na seção ``Referências Bibliográfica'', e devem ser citadas no texto consecutivamente em ordem numérica e entre colchetes.
% Exemplo: [1]
Pode-se gerar as referências bibliográficas através do comando \verb|\cite| (\cite{Lam94,ClarkAF}). O formato da bibliografia é determinado pelo arquivo de estilo bibliográfico ``aes.bst'' distribuído com esta classe.
%
\subsection{Estilo do texto}
O texto deve ser de fácil entendimento e gramaticamente correto. Palavras e frases não devem ser abreviadas em títulos e resumos, nem da primeira vez que aparecem.

Devem-se usar unidades métricas de acordo com o Sistema Internacional de Unidades (SI) \cite{Inmetro12}. Na Tabela~\ref{tab:units} estão listadas algumas unidades do SI mais freqüentemente usadas, e algumas que não são do SI.

\begin{table}[!ht]
\caption{Unidades SI e outras}
\label{tab:units}
\vspace*{10pt}
\centering
\small
\begin{tabular}{ll}
\textit{Nome da Unidade}      	&	\textit{Símbolo da Unidade}\\ \hline
ampere             		&	A\\
bit ou bits        		&	como escrito\\
bytes              		&	como escrito\\
decibel           		&	dB\\
ângulo (geométrico)     	&	$^\textrm{o}$\\
farad             		&	F\\
gauss             		&	Gs\\
grama              		&	g\\
henry              		&	H\\
hertz              		&	Hz\\
hora              		&	h\\
polegada             		&	in\\
joule              		&	J\\
kelvin             		&	K\\
kilohertz          		&	kHz\\
kilohm                		&	k$\Omega$\\
litro             		&	l, L\\
megahertz         		&	MHz\\
metro              		&	m\\
microfarad         		&	$\mu$F\\
micrometro         		&	$\mu$m\\
microsegundo        		&	$\mu$s\\
milliampere        		&	mA\\
millihenry         		&	mH\\
millimetro         		&	mm\\
millivolt          		&	mV\\
minuto (tempo)        		&	min\\
minuto (geométrico)     	&      	'\\
nanosegundo         		&	ns\\
oersted           		&	Oe\\
ohm                		&	$\Omega$\\
pascal             		&	Pa\\
picofarad          		&	pF\\
segundo (tempo)      		&	s\\
segundo (geométrico)    	&	"\\
siemens            		&	S\\
tesla              		&	T\\
volt               		&	V\\
watt               		&	W\\
weber              		&	Wb\\
\end{tabular} 
\end{table}
%
\section{DIREITO AUTORAL (COPYRIGHT)}
O texto entre linhas, contido no topo da primeira página do artigo de Convenção da AES Brasil, é de propriedade da \emph{Audio Engineering Society} e não pode ser reproduzido sem permissão. Os direitos sobre o conteúdo de um artigo de Convenção/Congresso da AES Brasil são do autor ou autores.

No entanto, submetendo um artigo para apresentação em uma Convenção/Congresso da AES, o autor estará concordando que o \emph{AES Journal} terá a preferência para publicação. Caso aceito para publicação no \emph{AES Journal} ou outro ``\emph{Special Issue}'' da \emph{AES}, será solicitada ao(s) autor(es) a transferência dos direitos autorais.
%

\section{FIGURAS E TABELAS}

Figuras, diagramas, gráficos, etc. têm que ser visíveis em impressão preto e branco (P\&B). Se forem coloridos, certifique-se de que as identificações serão coerentes em impressão P\&B. Isso pode ser feito com indicações textuais no próprio gráfico. Use linhas não-menores que 1/2 ponto. As figuras podem ocupar uma ou duas colunas.

Figuras, tabelas e ilustrações devem ser colocadas seqüencialmente no texto, e próximas do local onde se faz referência a elas, sempre que possível. Todas as figuras devem ser numeradas e referenciadas no texto por extenso (por exemplo: ``Figura~\ref{fig:logo}'', e não ``Fig.~\ref{fig:logo}'').

\begin{figure}[h!]
  \centering
	\ifpdfout
	    \includegraphics[width=2cm]{aeslogo.pdf}% (pdflatex)
	\else
		\includegraphics[width=2cm]{aeslogo.eps}% (latex)
	\fi
  \caption{Logotipo da AES}
  \label{fig:logo}
\end{figure}

Fotografias e imagens gráficas devem ser salvas em baixa resolução sempre que possível (de 72 a 300 dpi) desde que preservada a qualidade e legibilidade.

\section{EQUAÇÕES}
As equações devem ser numeradas seqüencialmente, entre parênteses, estando destacadas em linha própria. São citadas no texto da seguinte forma: ``Equação~(\ref{eq:1})''.

Para facilitar a composição de fórmulas, use-se a opção de classe \verb|ams|. Esta opção carrega os pacotes \verb|amsmath|, \verb|amssymb| e \verb|amsthm|. As equações podem ser destadacas em linha própria com o ambiente \verb|equation|. Por exemplo:
\begin{equation}
\label{eq:1}
  \left\{\,
    x\biggm|\int_{0}^x t^{2}\,dt\leq5
  \,\right\}.
\end{equation}

Equações que não couberem em uma coluna podem ser ajustadas em duas linhas; pode-se fazer isso com o pacote \textsf{breqn} usado para quebrar a equação em duas linhas.

%%%%%%%%%%%%%%%%%%%%%%%%%%%%%%%%%%%%%%%%%%%%%%%%%%%%%%%%%%%
%
% BIBLIOGRAPHY EXAMPLE 
%
% para gerar a bibliografia execute "bibtex template"
% to generate the bibliography please execute "bibtex template"
\bibliographystyle{aes} % style aes.bst
\bibliography{bib} % bibliography file in bibtex format
%
\end{document}


\end{document}